\section*{Chapter \ref{ch:b2:open-systems} --- Thermodynamic Balances and Open Systems}

\begin{solution}{exo:b2:open-systems:1}
(a) $\qty{2e-4}{m^3/s}/\qty{7.85e-5}{m^2} = \qty{2.5}{m/s}$. (b) $(0.1 \times 60 + 0.2
\times 10)/0.3 = \qty{26.7}{\celsius}$. (c) $\dot S_{\text{c}} = 0.1c\ln(299.8/333) +
0.2c\ln(299.8/283) = \qty{4.3}{W/K}$: heat passed from hot to cold water.
\end{solution}

\begin{solution}{exo:b2:open-systems:2}
$\sqrt{2 \times 200 \times 10^3 + 50^2} = \qty{634}{m/s}$; $c^2/2c_p = \qty{45}{K}$;
\qty{31}{K}.
\end{solution}

\begin{solution}{exo:b2:open-systems:3}
(a) $T$ unchanged, $s$ rises by $r\ln(P_1/P_2)$. (b) $h = \qty{2576}{kJ/kg}$:
$x = (2576 - 417)/2258 = 0.956$ --- the quality is read from the
temperature of the superheated (or just dry) outlet. (c) $(256 -
200)/199 = 0.28$.
\end{solution}

\begin{solution}{exo:b2:open-systems:4}
(a) \qty{41.8}{kW}; $41.8 \times 10^3/(1005 \times 20) = \qty{2.1}{kg/s}$, \qty{1.7}{m^3/s}.
(b) \qty{10}{MW}. (c) $v\Delta P = \qty{5.9}{kJ/kg}$; \qty{3}{MW}.
\end{solution}

\begin{solution}{exo:b2:open-systems:5}
(a) \cref{thm:b2:open-systems:firstlaw}. (b) A flow is pushed in and
out by its neighbours: that pressure work, $Pv$ per unit mass, is
bundled with $u$. (c) \qty{50}{J/kg} and \qty{45}{kJ/kg} against \qty{1}{kJ/kg}.
(d) \qty{981}{J/kg} against \qty{4180}{J/kg}: \qty{0.23}{K}.
\end{solution}

\begin{solution}{exo:b2:open-systems:6}
(a) \qty{531}{K}, \qty{239}{kJ/kg}. (b) \qty{299}{kJ/kg}, \qty{590}{K}, $c_p\ln(590/531)
= \qty{106}{J/kg/K}$. (c) Each stage ends at \qty{393}{K}: $2 \times 1005 \times 100 =
\qty{201}{kJ/kg}$, $16\%$ saved. (d) $rT\ln 8 = \qty{175}{kJ/kg}$; with
intercooling between many stages the path hugs the isotherm.
\end{solution}

\begin{solution}{exo:b2:open-systems:7}
(a) $x = (6.88 - 0.476)/7.918 = 0.81$, $h_{4s} = \qty{2098}{kJ/kg}$, $w_{\text{s}} =
\qty{1324}{kJ/kg}$. (b) \qty{1125}{kJ/kg}; $h_4 = 2297$, $x_4 = 0.89$. (c) $s_4 =
\qty{7.53}{kJ/kg/K}$: $s_{\text{c}} = \qty{0.65}{kJ/kg/K}$; lost $T_0s_{\text{c}} =
\qty{200}{kJ/kg}$. (d) Droplets at hundreds of metres per second erode the
last blades.
\end{solution}

\begin{solution}{exo:b2:open-systems:8}
(a) $h_2 = 399 + 27/0.8 = \qty{433}{kJ/kg}$; $w = \qty{34}{kJ/kg}$. (b) $q_{\text{evap}}
= 143$, $q_{\text{cond}} = \qty{177}{kJ/kg}$; $\mathrm{COP}_{\text{c}} = 4.2$,
$\mathrm{COP}_{\text{h}} = 5.2$. (c) $6.8$ and $7.8$; two thirds. (d) \qty{0.045}{kg/s},
\qty{1.5}{kW}; a colder evaporator means a lower pressure, a larger
pressure ratio (more work per kilogram) and a lighter vapour (less
mass per stroke).
\end{solution}

\begin{solution}{exo:b2:open-systems:9}
(a) \qty{30}{\celsius}. (b) $4180[\ln(323/363) + 2\ln(303/283)] = \qty{83}{W/K}$.
(c) \qty{24}{kW} of \qty{167}{kW}. (d) In parallel flow the temperature gaps are
larger and the hot outlet cannot fall below the cold outlet: more
entropy for the same duty.
\end{solution}

\begin{solution}{exo:b2:open-systems:10}
(a) $\dd(mu)/\dd t = \dot m h_0$ integrates to $mu = mh_0$: $c_vT = c_pT_0$,
$T = \gamma T_0$. (b) \qty{410}{K}. (c) The final gas is a mixture of the
original gas at $T_0$ and of incoming gas heated to $\gamma T_0$, hence
in between. (d) The flow work heats the gas; a slow fill in a water
bath lets the heat out so that the pressure at \qty{20}{\celsius} is the
one wanted.
\end{solution}

\begin{solution}{exo:b2:open-systems:11}
(a) $x = (7.76 - 0.476)/7.918 = 0.92$, $h = \qty{2367}{kJ/kg}$. (b) $w = 572 + 1111
- 6 = \qty{1677}{kJ/kg}$, $q = 3278 + 628 = \qty{3906}{kJ/kg}$, $\eta = 0.43$
against $0.40$. (c) The reheat adds heat at a high temperature and
raises the mean temperature of heat addition. (d) The exhaust is drier
($0.92$ instead of $0.81$).
\end{solution}

\begin{solution}{exo:b2:open-systems:12}
(a) \qty{293}{K}, \qty{596}{K}, \qty{1500}{K}, \qty{737}{K}. (b) $w_{\text{c}} = 305$,
$w_{\text{t}} = 767$, $q = \qty{909}{kJ/kg}$, $\eta = 0.51 = 1 - 12^{-0.286}$. (c)
Exhaust heat above \qty{373}{K}: \qty{366}{kJ/kg}; $0.7 \times 0.35 \times 366 =
\qty{90}{kJ/kg}$: $\eta = (462 + 90)/909 = 0.61$. (d) It takes heat at
\qty{1500}{K} and rejects it at \qty{306}{K}: the widest span any working
fluid allows.
\end{solution}

\begin{solution}{pb:b2:open-systems:1}
\textbf{1.} Compressor (adiabatic, $w = \Delta h$), condenser (isobaric, heat
out), \omterm{prop:b2:open-systems:devices}{throttle} (isenthalpic), evaporator (isobaric, heat in).

\textbf{2.} $h_4 = \qty{256}{kJ/kg}$, $x_4 = 0.28$.

\textbf{3.} \qty{143}{kJ/kg}.

\textbf{4.} \qty{27}{} and \qty{34}{kJ/kg}; $h_2 = 433$; $T_2 \approx 45 + 7 \approx
\qty{52}{\celsius}$.

\textbf{5.} \qty{177}{kJ/kg}; $\mathrm{COP}_{\text{h}} = 5.2$, $\mathrm{COP}_{\text{c}} = 4.2$;
$5.2 = 4.2 + 1$.

\textbf{6.} $7.8$ and $6.8$; two thirds.

\textbf{7.} \qty{0.045}{kg/s}; \qty{1.5}{kW}, \qty{1.7}{kW} at the plug.

\textbf{8.} Compressor: $c_p\ln(325/318) \approx \qty{0.022}{kJ/kg/K}$;
\omterm{prop:b2:open-systems:devices}{throttle}: $s_4 - s_3 = 1.204 - 1.19 = \qty{0.014}{kJ/kg/K}$ --- the compressor.

\textbf{9.} $313/55 = 5.7$; the pressure ratio nearly doubles (more work
per kilogram) and the vapour density halves (less mass per stroke):
capacity drops just when the house needs more.

\textbf{10.} The coil is below \qty{0}{\celsius} in humid air; frost
insulates it; the defrost reverses the cycle for minutes each hour and
costs a few per cent.

\textbf{11.} $308/35 = 8.8$ against $328/55 = 6.0$: the lower the
condenser temperature, the better --- underfloor.

\textbf{12.} \qty{1}{kWh} of fuel $\to$ \qty{0.4}{kWh} of electricity $\to$
\qty{1.6}{kWh} of heat, against \qty{0.9}{kWh} from the boiler.

\textbf{13.} Pump, boiler (isobaric heating, vaporisation, superheat),
turbine, condenser.

\textbf{14.} \qty{6}{kJ/kg}; $h_2 = 144$.

\textbf{15.} $3278 = 1069 + 1571 + 638$ (\unit{kJ/kg}).

\textbf{16.} $x = 0.81$, $h_{4s} = 2098$, $w_{\text{s}} = \qty{1324}{kJ/kg}$.

\textbf{17.} \qty{1125}{kJ/kg}, $h_4 = 2297$, $x_4 = 0.89$; wetter steam erodes the
blades and lowers the efficiency.

\textbf{18.} $\eta = 1119/3278 = 0.34$; Carnot $0.60$; $\bar T = 3278/6.40 =
\qty{512}{K}$, giving $0.40$.

\textbf{19.} \qty{536}{kg/s}; \qty{1160}{MW}; \qty{28}{m^3/s} of river water,
or \qty{480}{kg/s} evaporated.

\textbf{20.} \qty{0.65}{kJ/kg/K}, \qty{350}{kW/K}; \qty{200}{kJ/kg}, $18\%$ of the
turbine work.

\textbf{21.} The second heat addition is at high temperature (higher
$\bar T$), and the second expansion starts superheated and ends drier.

\textbf{22.} $1757\,\unit{MW} \times (1/512 - 1/1500) = \qty{2.3}{MW/K}$ --- six times
the turbine's.

\textbf{23.} $1160\,\unit{MW} \times (1/293 - 1/306) = \qty{0.17}{MW/K}$.

\textbf{24.} Furnace $\gg$ turbine $>$ condenser: hence higher steam
temperatures and pressures, reheat and regeneration; then better
blades; the condenser comes last.

\textbf{25.} Draw the \omterm{def:b2:open-systems:cv}{control volume}; mass in equals mass out; $\Delta h +
\Delta c^2/2 + g\Delta z = w_{\text{u}} + q$ with $h$ from $c_p$, tables or
$v\Delta P$; $\Delta s = \sum q/T + s_{\text{c}}$ with isentropic references and
efficiencies; assemble the components and compare with Carnot.
\end{solution}
